\section{Related Work}
\label{sec:related_work}

\subsection{Constellation Design and Optimization}

Parametric constellations significantly reduce the number of design parameters, with the Walker Delta pattern \citep{walker_satellite_1984} requiring only the number of planes, satellites per plane, and inclination. Flower constellations extend these ideas by using elliptical orbits to form specific ``harmonics'' in revisit rates to produce repeating ground tracks \citep{mortari_flower_2004}, allowing for directed coverage over a specific area. These parametric families are efficient but inherently symmetric, and cannot be tailored to non-uniform or mission-specific coverage requirements without breaking the symmetry assumptions that make them tractable.

For non-parametric design, previous work has relied primarily on metaheuristic methods.
\citet{paek_optimization_2019} uses genetic algorithms and simulated annealing to reconfigure constellations, optimizing for global coverage and regional revisit over five design variables, requiring approximately 2500 total evaluations through computationally expensive STK orbit simulations.
They also attempted gradient-based optimization via finite differences and found it ``extremely poorly conditioned,'' with convergence rarely achieved.
In a survey of 144 papers on constellation design, \citet{choo_survey_2024} finds that gradient methods appear as a design approach in only two of the surveyed articles, making it among the least common methods---ranking far behind genetic algorithms, simulated annealing, and analytical approaches.
Neither of the two entries actually uses analytic gradients.
The first, \citet{fraire_sparse_2020}, is labeled a gradient method but in fact uses a greedy neighbor search over discrete parameters (number of satellites, planes, inclination) with no gradients computed.
The second, \citet{abdelkhalik_optimization_2011}, pairs a genetic algorithm with a ``second-order gradient'' refinement step to cover ground sites under $J_2$ perturbations and explicitly notes that the fitness landscape contains numerous local minima that cause classical optimization methods to fail---but the gradient refinement is applied to a single-orbit subproblem rather than the multi-satellite design variables, so the overall design loop remains heuristic.
Separately, \citet{williams_rogers_optimal_2026} formulate constellation configuration as a collection of mixed-integer linear programs that select which satellites to activate from a pre-enumerated set of orbital slots, yielding provably optimal solutions for coverage and revisit objectives; however, the orbital slots and their underlying geometry (altitude, inclination, eccentricity) must be fixed a priori, so the method cannot discover improved orbit parameters.
We benchmark against SA/GA/DE baselines quantitatively in \secref{sec:baselines}.

\citet{hwang_large-scale_2014} demonstrates true gradient-based optimization of a satellite system through OpenMDAO (Open-Source Framework for Multidisciplinary Design, Analysis, and Optimization), optimizing a single satellite's subsystem design across seven coupled disciplines with over 25{,}000 design variables, with gradient computation made tractable with adjoint state methods. Their work utilizes gradient-based objectives combined with orbit-dependent parameters,  such as solar panel angles, battery sizing, and communication scheduling, to optimize a single spacecraft's operation. 

% This work differs in two respects: we use true analytical gradients through a complete computational graph rather than finite-difference approximations through a simulator, and we make the coverage and revisit objectives themselves differentiable via continuous relaxations rather than treating them as discrete metrics.

\subsection{Differentiable Physics and Simulation}

Differentiable programming~\citep{baydin_automatic_2018} allows for gradients to be obtained across domains even where forward models are complex but smooth. 
The benefits of differentiable programming have been made significantly more accessible through frameworks that perform automatic differentiation, such as PyTorch and JAX.
\citet{hu_difftaichi_2020} demonstrated differentiable physical simulators targeting robotics and fluid dynamics, and \citet{degrave_differentiable_2019} develops differentiable physics engines for contact-rich robotics.
For continuous relaxations applied to parameterized categorical distributions, Gumbel-Softmax reparameterization~\citep{jang_categorical_2017} is frequently used to compute analytic gradients.
\citet{kamra_gradient-based_2021} proves a coverage gradient theorem for spatial multi-resource coverage objectives, providing an estimator for the gradient that can be used for optimization of sensor placement, via spatial discretization and implicit boundary differentiation applied to 2D planar sensor networks.

In astrodynamics, \citet{acciarini_closing_2025} introduces $\partial$SGP4, a PyTorch reimplementation of the SGP4/SDP4 orbit propagation algorithm, which is forward mode differentiable.
\citet{naylor_dlite_2025} extends this direction by combining a JAX-based differentiable SGP4 propagator with Mitsuba~\citep{noauthor_mitsuba_nodate}, a differentiable renderer, to optimize trajectories for spacecraft-to-spacecraft inspection operations, demonstrating end-to-end forward-mode differentiation through orbit propagation chained with a rendering pipeline.

This work builds on $\partial$SGP4 by introducing differentiable coverage and revisit objectives that enable direct optimization of mission performance metrics.
